\documentclass[11pt]{article}
\usepackage[a4paper,margin=1in]{geometry}
\usepackage{amsmath,amssymb,amsfonts}
\usepackage{mathtools}
\usepackage{hyperref}
\usepackage{enumitem}
\usepackage[title]{appendix}
\hypersetup{colorlinks=true,linkcolor=blue,citecolor=blue,urlcolor=blue}
\setlist{nosep}
\numberwithin{equation}{section}

\title{Unified Boundary Field Theory of Interest Rates III: The Stochastic Long Boundary \\\large Volatility-Aligned Geometry, Affine Admissibility, and Two-Channel MSNR for Instantaneous Rolling Forward Rates}

\author{Sachin Oza}
\date{\today}

\begin{document}
\maketitle

\begin{abstract}
This paper completes the theoretical construction of the Unified Boundary Field Theory for instantaneous rolling forward rates (IRFRs). Earlier parts of the theory show that a one-boundary forward-rate system admits non-unique affine drift and volatility structures with constant market signal-to-noise ratios. The present paper extends that construction by adding a stochastic long-boundary and deriving the corresponding two-boundary geometry.

The central object is the maturity-loading relation
\[
x_{t,\tau+t}-X_t
=
\left(x_{t,t}-X_t\right)e^{-a^2\tau},
\]
which expresses the remaining-maturity \(\tau\) instantaneous rolling forward rate, at calendar time \(t\), as an exponentially weighted interpolation between the short-boundary state, \(x_{t,t}\), and the stochastic long-boundary, \(X_t\). This relation provides both an expected shape condition and a local volatility-alignment condition.

The resulting two-boundary geometry identifies an affine-preserving free-boundary subclass in which volatility alignment, drift-to-volatility proportionality, and a single structural Brownian correlation are mutually compatible. This is an existence and preservation result for a tractable subclass, not a classification theorem for all possible two-boundary transformations. Within this tractable subclass, the completed two-boundary system has a constant market signal-to-noise ratio (MSNR). The Robin branch is then reformulated as a boundary-layer coordinate representation: a partially anchored endpoint, its local maturity slope, and the long boundary are projected through a deterministic loading vector. The MSNR condition is interpreted as a local return-to-risk discipline on this projection. A rank-two Robin covariance closure is then isolated as a natural stochastic special case: the three Robin boundary-layer correlations lie on a two-dimensional covariance manifold and one pairwise correlation is determined, up to orientation, by the other two. In the symmetric affine-preserving closure, the same parameters also define a coherent two-boundary equilibrium density: a correlated inverse-square-root extension of the shifted inverse-gamma law that appears in the one-boundary equilibrium case. Along the matched-boundary diagonal of the two-boundary density, the exponential concentration rate is exactly the symmetric MSNR. The result provides the analytic bridge between the one-boundary admissibility theory, the symmetric two-boundary equilibrium law, and the empirical two-boundary covariance structure.
\end{abstract}

\section{Introduction}

The instantaneous rolling forward rate (IRFR) boundary-field framework begins from the premise that the term structure is a maturity field constrained by boundary conditions. Paper 1 developed the basic maturity-field representation and the exponential boundary transport clock \cite{OzaIRFR}. Paper 2 classified short end boundary regimes and showed that the Robin or mixed short end condition generates the repeated root loading \cite{OzaShortBoundary}
\[
 \tau e^{-a^2\tau}.
\]
This paper completes the theoretical construction by adding a stochastic realised long-maturity boundary.

The fixed long-boundary theory identifies the asymptotic state with an equilibrium anchor \(x_\infty\). That is useful as a long-run reference, but it is too restrictive when the realised long end of the curve can move persistently. The present paper therefore separates the equilibrium anchor \(x_\infty\) from the realised long-boundary state \(X_t\). The long-boundary displacement is
\begin{equation}
\label{eq:Y-definition}
Y_t:=X_t-x_\infty.
\end{equation}

The completed maturity geometry is generated by three loadings:
\begin{equation}
\label{eq:three-loadings}
e^{-a^2\tau},
\qquad
\tau e^{-a^2\tau},
\qquad
1-e^{-a^2\tau}.
\end{equation}
\begin{itemize}
    \item The first loading transports the realised short boundary into maturity space. 
    \item The second is the Robin interior deformation inherited from the short-boundary classification. 
    \item The third transports the stochastic long-boundary into maturity space.
\end{itemize}

In the canonical bridge representation the curve is written
\begin{equation}
\label{eq:canonical-bridge-intro}
x_t(\tau)
=
x_{t,t}e^{-a^2\tau}
+
B_tx_{t,t}\tau e^{-a^2\tau}
+
X_t\left(1-e^{-a^2\tau}\right).
\end{equation}
The corresponding maturity-space Jacobian is
\begin{equation}
\label{eq:intro-jacobian}
J(\tau)
=
\left(
 e^{-a^2\tau},
 \tau e^{-a^2\tau},
 1-e^{-a^2\tau}
\right).
\end{equation}
The first two components are inherited from the short end boundary theory. The third is the new stochastic long-boundary component.

The main theoretical question is whether this added long-boundary preserves an affine admissible subclass of the one-boundary construction, and the answer is yes. The long-boundary imposes a compatibility condition on the coordinate geometry, but it does not remove the affine-preserving subclass identified below. The required compatibility is volatility alignment:
\begin{equation}
\label{eq:intro-vol-align}
y_x=\frac{g_y}{g_x},
\qquad
Y_X=\frac{g_Y}{g_X}.
\end{equation}
Once this is imposed, the remaining diffusion equations identify the compatible correlation branches \(\rho_{yY}=\pm\rho_{xX}\). On the orientation-preserving branch, affine \(f\) and \(g\) functions, constant \(f/g\), and a single constant correlation preserve a global two-boundary market signal to noise ratio (MSNR).
The paper proceeds as follows:

\begin {itemize}
 \item Section 2 states the stochastic long-boundary decomposition. \item Section 3 records the instantaneous expected level and full instantaneous volatility implied by the exponential maturity loading. \item Section 4 derives the two-variable Kolmogorov forward equation. \item Section 5 performs the coordinate transform and coefficient matching. \item Section 6 states the affine-preserving admissible subclass and derives the corresponding constant correlated MSNR. \item Section 7 reformulates the Robin branch as a boundary-layer coordinate representation, derives the projected covariance formula, and gives a dictionary separating structural, diagnostic, and reduced correlation objects. \item Section 8 isolates a rank-two Robin covariance closure and the associated boundary-correlation triangle. \item Section 9 gives the symmetric two-boundary equilibrium density associated with the compact closure. \item Section 10 records the asymptotic boundary dynamics and long-end variance floor. \item Sections 11--13 discuss maturity-space rank, risk-neutral compatibility, empirical implications, and the summary of the construction.
\end {itemize}


\section{Stochastic long-boundary decomposition}

Let \(\tau=T-t\) denote remaining maturity. In the fixed long-boundary theory, the asymptotic boundary is represented by an equilibrium anchor \(x_\infty\). The stochastic long-boundary extension separates this equilibrium anchor from the realised long-boundary state \(X_t\). We write the realised displacement from equilibrium as
\[
Y_t:=X_t-x_\infty.
\]

The long-boundary loading is
\[
1-e^{-a^2\tau}.
\]
It vanishes at the short end and tends to one at the long end. Thus the stochastic long-boundary contribution is
\[
\left(1-e^{-a^2\tau}\right)(X_t-x_\infty).
\]

For a short-end boundary branch \(j\) (which, in paper 2, classifies Neumann/free, Dirichlet/fixed and Robin/mixed branches),  the maturity field may be decomposed as
\begin{equation}
\label{eq:boundary-decomposition}
x_t^j(\tau)
=
x_\infty\left(1-e^{-a^2\tau}\right)
+
\left(1-e^{-a^2\tau}\right)(X_t-x_\infty)
+
Z_t^j(\tau).
\end{equation}
The first term on the RHS is the fixed-anchor transport component. The second term transports the realised long-boundary displacement into maturity space. The residual \(Z_t^j(\tau)\) is the finite-maturity shape component associated with the short-end boundary branch.

Equivalently, since \(Y_t=X_t-x_\infty\),
\[
x_t^j(\tau)
=
x_\infty\left(1-e^{-a^2\tau}\right)
+
\left(1-e^{-a^2\tau}\right)Y_t
+
Z_t^j(\tau).
\]
The split form is retained because it distinguishes the equilibrium anchor from the realised long-boundary displacement, even though the first two terms combine algebraically to
\[
X_t\left(1-e^{-a^2\tau}\right).
\]

At the short end,
\[
1-e^{-a^2\cdot 0}=0,
\]
so the stochastic long boundary does not move the instantaneous endpoint. At the long end,
\[
1-e^{-a^2\tau}\to1
\qquad
\text{as}
\qquad
\tau\to\infty.
\]
Therefore, if the finite-maturity residual satisfies
\[
Z_t^j(\tau)\to0,
\]
then
\[
\lim_{\tau\to\infty}x_t^j(\tau)=X_t.
\]
The stochastic long-boundary extension therefore replaces the fixed asymptotic limit \(x_\infty\) with the realised long-boundary state \(X_t\), while preserving \(x_\infty\) as the equilibrium reference point.

For the Robin branch, the canonical bridge is
\begin{equation}
\label{eq:canonical-bridge}
x_t(\tau)
=
x_{t,t}e^{-a_R^2\tau}
+
B_tx_{t,t}\tau e^{-a_R^2\tau}
+
X_t\left(1-e^{-a_R^2\tau}\right).
\end{equation}
The short endpoint is pinned to \(x_{t,t}\), the long endpoint is pinned to \(X_t\), and the interior deformation is carried by the repeated-root Robin loading
\[
\tau e^{-a_R^2\tau},
\]
which vanishes at both endpoints. The parameter \(a_R\) is written separately to allow the Robin maturity clock to differ from the corresponding Neumann/free (and indeed the Dirichlet/fixed) clock.

\section{Instantaneous expected level and volatility loading}

The original maturity-loading intuition is most transparent when \(T_x=t\). Then the short-side coordinate is \(x_{t,t}\), the maturity-\(T\) coordinate is \(x_{t,T}\), and the long-boundary state is \(X_t\). The exponential boundary relation is
\begin{equation}
\label{eq:expected-level-boundary-relation}
x_{t,T}-X_t
=
\left(x_{t,t}-X_t\right)e^{-a^2(T-t)}.
\end{equation}
Equivalently,
\begin{equation}
\label{eq:expected-level-direct}
x_{t,T}
=
e^{-a^2(T-t)}x_{t,t}
+
\left(1-e^{-a^2(T-t)}\right)X_t.
\end{equation}
Therefore the instantaneous conditional expected level at maturity \(T\), given \(x_{t,t}\) and \(X_t\), is
\begin{equation}
\label{eq:instantaneous-expected-level}
\mathbb E_t[x_{t,T}]
=
e^{-a^2(T-t)}x_{t,t}
+
\left(1-e^{-a^2(T-t)}\right)X_t.
\end{equation}
Equivalently,
\[
\mathbb E_t[x_{t,T}-X_t]
=
\left(x_{t,t}-X_t\right)e^{-a^2(T-t)}.
\]

Differentiating \eqref{eq:expected-level-direct} with respect to \(x_{t,t}\), holding \(X_t\) fixed, gives
\begin{equation}
\label{eq:local-maturity-loading-derivative}
\frac{\partial x_{t,T}}{\partial x_{t,t}}
=
e^{-a^2(T-t)}.
\end{equation}
The coefficient-matching result below requires local volatility alignment:
\[
\frac{\partial x_{t,T}}{\partial x_{t,t}}
=
\frac{g_{t,T}}{g_{t,t}}.
\]
Therefore
\begin{equation}
\label{eq:maturity-vol-alignment}
\frac{g_{t,T}}{g_{t,t}}
=
e^{-a^2(T-t)},
\qquad
\text{or}
\qquad
g_{t,T}=e^{-a^2(T-t)}g_{t,t}.
\end{equation}
This is the local maturity-loading volatility relation when the long-boundary is held fixed.

If the long-boundary is itself stochastic, the full instantaneous volatility of \(x_{t,T}\) also contains boundary risk. Suppose
\begin{equation}
\label{eq:short-sde}
dx_{t,t}=-f_t\,dt+\sqrt{2}\,g_t\,dW_t^t,
\end{equation}
and
\begin{equation}
\label{eq:X-sde-simple}
dX_t=-f_X\,dt+\sqrt{2}\,g_X\,dW_t^X,
\end{equation}
with
\begin{equation}
\label{eq:rho-simple}
dW_t^t\,dW_t^X=\rho\,dt.
\end{equation}
Over an instantaneous step, treating the maturity-loading coefficient as fixed, \eqref{eq:expected-level-direct} gives
\[
dx_{t,T}
=
e^{-a^2(T-t)}dx_{t,t}
+
\left(1-e^{-a^2(T-t)}\right)dX_t.
\]
Substitution of \eqref{eq:short-sde} and \eqref{eq:X-sde-simple} yields
\begin{align}
\label{eq:maturity-T-sde}
dx_{t,T}
&=
-\left[
e^{-a^2(T-t)}f_t
+
\left(1-e^{-a^2(T-t)}\right)f_X
\right]dt
\nonumber\\
&\quad+
\sqrt{2}\left[
e^{-a^2(T-t)}g_t\,dW_t^t
+
\left(1-e^{-a^2(T-t)}\right)g_X\,dW_t^X
\right].
\end{align}
Hence the instantaneous drift coefficient is
\begin{equation}
\label{eq:instantaneous-drift-maturity-T}
f_{t,T}
=
e^{-a^2(T-t)}f_t
+
\left(1-e^{-a^2(T-t)}\right)f_X.
\end{equation}
The full two-boundary instantaneous variance coefficient is
\begin{align}
\label{eq:instantaneous-vol-maturity-T}
g_{t,T}^2
&=
e^{-2a^2(T-t)}g_t^2
+
\left(1-e^{-a^2(T-t)}\right)^2g_X^2
\nonumber\\
&\quad+
2\rho e^{-a^2(T-t)}\left(1-e^{-a^2(T-t)}\right)g_tg_X.
\end{align}
Thus the exponential maturity loading governs both the instantaneous expected level and the local volatility alignment, while the full stochastic-boundary volatility is a correlated combination of short-side and long-boundary risk.

\section{Two variable KFE before transformation}

We now record the two variable Kolmogorov forward equation that underlies the coefficient matching argument. Let the original coordinates be \((x,X)\), where \(x\) is the short-side or interior state and \(X\) is the long-boundary state. Using the conventions in papers 1 and 2,
\begin{equation}
\label{eq:xX-sde}
dx_t=-f_x\,dt+\sqrt{2}\,g_x\,dW_t^x,
\qquad
dX_t=-f_X\,dt+\sqrt{2}\,g_X\,dW_t^X,
\end{equation}
with
\begin{equation}
\label{eq:xX-corr}
dW_t^x\,dW_t^X=\rho_{xX}\,dt.
\end{equation}
For a density \(\varphi=\varphi(x,X,t)\), the KFE is
\begin{align}
\label{eq:xX-KFE-conservative}
\frac{\partial\varphi}{\partial t}
&=
\frac{\partial}{\partial x}(f_x\varphi)
+
\frac{\partial}{\partial X}(f_X\varphi)
+
\frac{\partial^2}{\partial x^2}(g_x^2\varphi)
\nonumber\\
&\quad+
2\rho_{xX}\frac{\partial^2}{\partial x\,\partial X}(g_xg_X\varphi)
+
\frac{\partial^2}{\partial X^2}(g_X^2\varphi).
\end{align}
Expanding and collecting derivatives of \(\varphi\), write
\begin{equation}
\label{eq:xX-KFE-expanded}
\frac{\partial\varphi}{\partial t}
=
A_0\varphi+A_x\varphi_x+A_X\varphi_X+A_{xx}\varphi_{xx}+A_{xX}\varphi_{xX}+A_{XX}\varphi_{XX},
\end{equation}
where
\begin{align}
\label{eq:A0}
A_0
&=
\frac{\partial f_x}{\partial x}
+
\frac{\partial f_X}{\partial X}
+
\frac{\partial^2 g_x^2}{\partial x^2}
+
2\rho_{xX}\frac{\partial^2(g_xg_X)}{\partial x\,\partial X}
+
\frac{\partial^2 g_X^2}{\partial X^2},
\end{align}
\begin{align}
\label{eq:AxAX}
A_x
&=
f_x+2\frac{\partial g_x^2}{\partial x}
+2\rho_{xX}\frac{\partial(g_xg_X)}{\partial X},
\\
A_X
&=
f_X+2\frac{\partial g_X^2}{\partial X}
+2\rho_{xX}\frac{\partial(g_xg_X)}{\partial x},
\end{align}
and
\begin{equation}
\label{eq:A-second}
A_{xx}=g_x^2,
\qquad
A_{xX}=2\rho_{xX}g_xg_X,
\qquad
A_{XX}=g_X^2.
\end{equation}

\section{Coordinate transform and coefficient matching}

Transform from \((x,X)\) to \((y,Y)\), where
\begin{equation}
\label{eq:coordinate-transform}
y=y(x,X),
\qquad
Y=Y(X).
\end{equation}
Since \(Y\) depends only on \(X\),
\[
\frac{dY}{dx}:=Y_x=0,
\qquad
Y_X=\frac{dY}{dX},
\qquad
Y_{XX}=\frac{d^2Y}{dX^2}.
\]
For \(\varphi=\varphi(y,Y,t)\), the chain-rule identities are
\begin{equation}
\label{eq:first-derivatives}
\varphi_x=y_x\varphi_y,
\qquad
\varphi_X=y_X\varphi_y+Y_X\varphi_Y,
\end{equation}
\begin{equation}
\label{eq:second-xx}
\varphi_{xx}=y_x^2\varphi_{yy}+y_{xx}\varphi_y,
\end{equation}
\begin{equation}
\label{eq:second-xX}
\varphi_{xX}
=
y_xy_X\varphi_{yy}
+
y_xY_X\varphi_{yY}
+
y_{xX}\varphi_y,
\end{equation}
and
\begin{align}
\label{eq:second-XX}
\varphi_{XX}
&=
y_X^2\varphi_{yy}
+
2y_XY_X\varphi_{yY}
+
Y_X^2\varphi_{YY}
+
y_{XX}\varphi_y
+
Y_{XX}\varphi_Y.
\end{align}
Substituting \eqref{eq:first-derivatives}--\eqref{eq:second-XX} into \eqref{eq:xX-KFE-expanded} gives
\begin{equation}
\label{eq:yY-KFE-B}
\frac{\partial\varphi}{\partial t}
=
B_0\varphi+B_y\varphi_y+B_Y\varphi_Y+B_{yy}\varphi_{yy}+B_{yY}\varphi_{yY}+B_{YY}\varphi_{YY},
\end{equation}
with
\begin{equation}
\label{eq:B0}
B_0=A_0,
\end{equation}
\begin{equation}
\label{eq:By}
B_y
=
A_xy_x+A_Xy_X+A_{xx}y_{xx}+A_{xX}y_{xX}+A_{XX}y_{XX},
\end{equation}
\begin{equation}
\label{eq:BY}
B_Y=A_XY_X+A_{XX}Y_{XX},
\end{equation}
\begin{equation}
\label{eq:Byy}
B_{yy}
=
A_{xx}y_x^2+A_{xX}y_xy_X+A_{XX}y_X^2,
\end{equation}
\begin{equation}
\label{eq:ByY}
B_{yY}=A_{xX}y_xY_X+2A_{XX}y_XY_X,
\end{equation}
and
\begin{equation}
\label{eq:BYY}
B_{YY}=A_{XX}Y_X^2.
\end{equation}

In the transformed coordinates, the diffusion-side coefficients have the same structural form:
\begin{equation}
\label{eq:yY-diffusion-coefficients}
B_{yy}=g_y^2,
\qquad
B_{yY}=2\rho_{yY}g_yg_Y,
\qquad
B_{YY}=g_Y^2.
\end{equation}
Therefore the diffusion matching equations are
\begin{equation}
\label{eq:diff-match-yy}
g_y^2
=
g_x^2y_x^2
+
2\rho_{xX}g_xg_Xy_xy_X
+
g_X^2y_X^2,
\end{equation}
\begin{equation}
\label{eq:diff-match-yY}
2\rho_{yY}g_yg_Y
=
2\rho_{xX}g_xg_Xy_xY_X
+
2g_X^2y_XY_X,
\end{equation}
and
\begin{equation}
\label{eq:diff-match-YY}
g_Y^2=g_X^2Y_X^2.
\end{equation}
The final equation gives the long-boundary volatility alignment
\begin{equation}
\label{eq:Y-vol-alignment}
Y_X=\frac{g_Y}{g_X},
\end{equation}
where the positive branch preserves the orientation of the long-boundary coordinate. We impose the analogous volatility-aligned orientation for the interior maturity coordinate,
\begin{equation}
\label{eq:y-vol-alignment}
y_x=\frac{g_y}{g_x}.
\end{equation}
Equations \eqref{eq:Y-vol-alignment} and \eqref{eq:y-vol-alignment} are the geometric compatibility conditions of the long-boundary transform. They are the coefficient-matching version of the maturity-loading identity \eqref{eq:maturity-vol-alignment}.

Once both volatility alignment relations are imposed, the two Brownian correlations are not independent. Substituting
\[
Y_X=\frac{g_Y}{g_X},
\qquad
 y_x=\frac{g_y}{g_x}
\]
into the mixed diffusion equation \eqref{eq:diff-match-yY} gives
\[
\rho_{yY}g_yg_Y
=
\rho_{xX}g_xg_X
\left(\frac{g_y}{g_x}\right)
\left(\frac{g_Y}{g_X}\right)
+
g_X^2y_X\left(\frac{g_Y}{g_X}\right).
\]
Therefore
\[
\rho_{yY}g_yg_Y=\rho_{xX}g_yg_Y+g_Xg_Yy_X,
\]
and, away from the degenerate case \(g_Y=0\),
\begin{equation}
\label{eq:yX-forced}
y_X
=
\frac{g_y}{g_X}\left(\rho_{yY}-\rho_{xX}\right).
\end{equation}
Substituting \eqref{eq:y-vol-alignment} and \eqref{eq:yX-forced} into the \(yy\) diffusion equation \eqref{eq:diff-match-yy} yields
\[
g_y^2
=
g_y^2
+
2\rho_{xX}g_y^2\left(\rho_{yY}-\rho_{xX}\right)
+
g_y^2\left(\rho_{yY}-\rho_{xX}\right)^2.
\]
Cancelling the first term on the right hand side and factoring gives
\begin{equation}
\label{eq:rho-branch-constraint}
g_y^2
\left(\rho_{xX}-\rho_{yY}\right)
\left(\rho_{xX}+\rho_{yY}\right)=0.
\end{equation}
Thus, away from the degenerate case \(g_y=0\), the compatible branches are
\begin{equation}
\label{eq:rho-branches}
\rho_{yY}=\rho_{xX}
\qquad\text{or}\qquad
\rho_{yY}=-\rho_{xX}.
\end{equation}
The first branch is the orientation-preserving branch: the coordinate transformation preserves the sign of the Brownian correlation. We therefore use the single structural correlation
\begin{equation}
\label{eq:single-rho-branch}
\rho_{yY}=\rho_{xX}=\rho
\end{equation}
in the admissible affine specification. The opposite branch corresponds to a sign reversal in one coordinate orientation rather than to an independent second correlation parameter.

The \(Y\)-drift coefficient gives the useful separation relation. From \eqref{eq:BY}, using \(Y_X=g_Y/g_X\) and \(A_{XX}=g_X^2\),
\begin{equation}
\label{eq:BY-substitute}
B_Y
=
A_X\frac{g_Y}{g_X}
+
g_X^2\frac{d}{dX}\left(\frac{g_Y}{g_X}\right).
\end{equation}
Substituting \(A_X\) gives
\begin{align}
\label{eq:BY-expanded}
B_Y
&=
\frac{g_Y}{g_X}
\left[
f_X+2\frac{\partial g_X^2}{\partial X}
+2\rho_{xX}\frac{\partial(g_xg_X)}{\partial x}
\right]
+
g_X^2\frac{d}{dX}\left(\frac{g_Y}{g_X}\right).
\end{align}
Since \(g_Y=g_Y(Y)\),
\[
\frac{dg_Y}{dX}
=
\frac{\partial g_Y}{\partial Y}\frac{dY}{dX}
=
\frac{\partial g_Y}{\partial Y}\frac{g_Y}{g_X}.
\]
Therefore
\begin{equation}
\label{eq:quotient-term}
g_X^2\frac{d}{dX}\left(\frac{g_Y}{g_X}\right)
=
g_Y\frac{\partial g_Y}{\partial Y}
-
g_Y\frac{\partial g_X}{\partial X}.
\end{equation}
The quotient rule and the independence assumptions remove the remaining mixed terms, yielding a clean separation between the long-boundary and short-side dynamics. Matching \(B_Y\) against the transformed coefficient
\[
B_Y=f_Y+2\frac{\partial g_Y^2}{\partial Y}+2\rho_{yY}\frac{\partial(g_yg_Y)}{\partial y},
\]
dividing by \(g_Y\), and simplifying gives
\begin{align}
\label{eq:separated-relation-general}
\frac{f_Y}{g_Y}
+
3\frac{\partial g_Y}{\partial Y}
+
\frac{2\rho_{yY}\partial_y(g_yg_Y)}{g_Y}
&=
\frac{f_X}{g_X}
+
3\frac{\partial g_X}{\partial X}
+
\frac{2\rho_{xX}\partial_x(g_xg_X)}{g_X}.
\end{align}
If \(g_Y\) is independent of \(y\), \(g_X\) is independent of \(x\), and the orientation-preserving branch \eqref{eq:single-rho-branch} is used, this reduces to
\begin{equation}
\label{eq:separated-relation-clean}
\frac{f_Y}{g_Y}
+
3\frac{\partial g_Y}{\partial Y}
+
2\rho\frac{\partial g_y}{\partial y}
=
\frac{f_X}{g_X}
+
3\frac{\partial g_X}{\partial X}
+
2\rho\frac{\partial g_x}{\partial x}.
\end{equation}
This is the coefficient-matching separation equation for the long-boundary transform.

\section{Affine admissibility and constant MSNR}

The affine relations used here are inherited from the one-boundary IRFR construction. The stochastic long-boundary extension does not require these relations to be unique. Rather, the paper restricts attention to the affine-preserving subclass: volatility-aligned coordinate maps for which the one-boundary affine drift--volatility closure survives after the long-boundary state is added. Within this subclass the two-boundary MSNR is explicit.

Thus the claim in this section is deliberately modest. Given a stochastic long boundary, there exists a tractable volatility-aligned affine subclass that preserves constant MSNR. The paper does not claim uniqueness of that subclass, nor does it classify every possible two-boundary coordinate transformation.

We now impose the affine admissible class. In the original coordinates, let
\begin{equation}
\label{eq:old-affine-g}
g_x(x)=p_x+q_xx,
\qquad
g_X(X)=p_X+q_XX,
\end{equation}
and
\begin{equation}
\label{eq:old-affine-f}
f_x=\gamma g_x,
\qquad
f_X=\delta g_X,
\end{equation}
where \(\gamma\) and \(\delta\) are constants. Thus
\[
q_x=\frac{\partial g_x}{\partial x},
\qquad
q_X=\frac{\partial g_X}{\partial X}.
\] Assume also that the orientation-preserving branch gives the single constant correlation
\[
\rho_{xX}=\rho_{yY}=\rho.
\]
Then
\[
\frac{f_x}{g_x}=\gamma,
\qquad
\frac{f_X}{g_X}=\delta.
\]
Thus the original-coordinate component MSNRs are constant.

In the transformed coordinates, impose the analogous affine structure
\begin{equation}
\label{eq:new-affine-g}
g_y(y)=p_y+q_yy,
\qquad
g_Y(Y)=p_Y+q_YY,
\end{equation}
and
\begin{equation}
\label{eq:new-affine-f}
f_y=\alpha g_y,
\qquad
f_Y=\beta g_Y,
\end{equation}
where \(\alpha\) and \(\beta\) are constants. Thus
\[
q_y=\frac{\partial g_y}{\partial y},
\qquad
q_Y=\frac{\partial g_Y}{\partial Y}.
\] The same \(\rho\) is used in the transformed coordinates by \eqref{eq:single-rho-branch}.
Then
\begin{equation}
\label{eq:component-MSNRs}
\frac{f_y}{g_y}=\alpha,
\qquad
\frac{f_Y}{g_Y}=\beta.
\end{equation}
The transformed-coordinate component MSNRs are therefore constant.

The clean separation relation \eqref{eq:separated-relation-clean} becomes
\begin{equation}
\label{eq:affine-separation}
\beta+3q_Y+2\rho q_y
=
\delta+3q_X+2\rho q_x.
\end{equation}
Equivalently,
\begin{equation}
\label{eq:affine-separation-zero}
(\beta-\delta)+3(q_Y-q_X)+2\rho(q_y-q_x)=0.
\end{equation}
Thus the affine family remains admissible subject to the algebraic compatibility condition \eqref{eq:affine-separation-zero}, the volatility-aligned geometric conditions \eqref{eq:Y-vol-alignment}--\eqref{eq:y-vol-alignment}, and the orientation-preserving correlation branch \eqref{eq:single-rho-branch}. The class is non-unique because many choices of affine coefficients and coordinate maps can satisfy these relations. The result should therefore be read as an existence and preservation result for a tractable affine subclass, not as a classification theorem for all volatility-aligned two-boundary transformations.

The correlated scalar MSNR is the quadratic form determined by the diffusion covariance. In the transformed coordinates, the diffusion matrix without the common factor \(2\) is
\begin{equation}
\label{eq:Sigma-yY}
\Sigma
=
\begin{pmatrix}
g_y^2 & \rho g_yg_Y \\
\rho g_yg_Y & g_Y^2
\end{pmatrix}.
\end{equation}
The squared MSNR is
\begin{equation}
\label{eq:MSNR-quadratic-form}
\mathrm{MSNR}_{y,Y}^2
=
\begin{pmatrix} f_y & f_Y \end{pmatrix}
\Sigma^{-1}
\begin{pmatrix} f_y \\ f_Y \end{pmatrix}.
\end{equation}
Since
\[
\det(\Sigma)=g_y^2g_Y^2(1-\rho^2),
\]
we have
\[
\Sigma^{-1}
=
\frac{1}{g_y^2g_Y^2(1-\rho^2)}
\begin{pmatrix}
g_Y^2 & -\rho g_yg_Y \\
-\rho g_yg_Y & g_y^2
\end{pmatrix}.
\]
Therefore
\begin{align}
\label{eq:MSNR-expanded}
\mathrm{MSNR}_{y,Y}^2
&=
\frac{1}{1-\rho^2}
\left[
\left(\frac{f_y}{g_y}\right)^2
-2\rho\left(\frac{f_y}{g_y}\right)\left(\frac{f_Y}{g_Y}\right)
+
\left(\frac{f_Y}{g_Y}\right)^2
\right].
\end{align}
Using \eqref{eq:component-MSNRs}, this becomes
\begin{equation}
\label{eq:constant-correlated-MSNR}
\boxed{
\mathrm{MSNR}_{y,Y}^2
=
\frac{\alpha^2-2\rho\alpha\beta+\beta^2}{1-\rho^2}.
}
\end{equation}
Hence
\begin{equation}
\label{eq:constant-correlated-MSNR-root}
\boxed{
\mathrm{MSNR}_{y,Y}
=
\left(\frac{\alpha^2-2\rho\alpha\beta+\beta^2}{1-\rho^2}\right)^{1/2}.
}
\end{equation}
Since \(\alpha\), \(\beta\), and \(\rho\) are constants, the scalar correlated MSNR is a global constant.

In the symmetric case \(\alpha=\beta=a\),
\[
\mathrm{MSNR}_{y,Y}^2=\mathrm{MSNR}^2
=
\frac{a^2-2\rho a^2+a^2}{1-\rho^2}
=
\frac{2a^2(1-\rho)}{(1-\rho)(1+\rho)}.
\]
For \(\rho\neq1\),
\begin{equation}
\label{eq:symmetric-MSNR}
\boxed{
\mathrm{MSNR}^2=\frac{2a^2}{1+\rho},
\qquad
\mathrm{MSNR}=|a|\sqrt{\frac{2}{1+\rho}}.
}
\end{equation}
This is the two-boundary analogue of the one-boundary constant MSNR clock.


\section{Robin boundary-layer coordinates and projected covariance}

The coefficient-matching construction in the preceding sections is naturally formulated in the free, or Neumann, boundary coordinates
\[
        Z^N_t=(x^N_{t,t},X_t)^\top,
\]
where \(x^N_{t,t}\) denotes the free short-end state and \(X_t\) denotes the stochastic long-boundary state. In that representation the instantaneous Brownian correlation \(\rho_N\) is the structural two-state KFE correlation.

The Robin case should not be obtained by simply replacing \(x^N_{t,t}\) with a fitted Robin endpoint in the same two-state KFE. The Robin endpoint is a boundary-transformed object: it is partially anchored between the free short boundary \(x^N_{t,t}\) and a pinned Dirichlet boundary \(x_D\). A useful geometric guide is
\begin{equation}
\label{eq:robin-neumann-dirichlet-heuristic}
        x^R_{t,t}
        \approx
        \omega_t x^N_{t,t}+(1-\omega_t)x_D,
        \qquad 0\leq \omega_t\leq 1.
\end{equation}
Here \(\omega_t=1\) corresponds to the Neumann limit and \(\omega_t=0\) corresponds to the Dirichlet limit. This interpolation is not imposed as a new theorem; it is a roadmap for the local coordinates used below. It emphasises that the Robin endpoint is neither the free short state nor the pinned state, but a partially anchored boundary value.

For a fixed local regime, write the Robin bridge as
\begin{equation}
\label{eq:robin-bridge-local}
        x^R_{t,t+\tau}
        =
        X_t+\left(A_t+B_t\tau\right)e^{-a_R^2\tau},
        \qquad \tau=T-t.
\end{equation}
At the short end,
\begin{equation}
\label{eq:robin-short-end-A}
        x^R_{t,t}=X_t+A_t,
\end{equation}
so
\begin{equation}
\label{eq:A-def-robin-local}
        A_t=x^R_{t,t}-X_t.
\end{equation}
Thus \(A_t\) is the Robin short-long boundary displacement.

The local short-end slope is defined by
\begin{equation}
\label{eq:s-def-robin-local}
        s_t
        =
        \left.
        \frac{\partial x^R_{t,t+\tau}}{\partial \tau}
        \right|_{\tau=0}.
\end{equation}
Differentiating the Robin bridge gives
\begin{equation}
\label{eq:s-B-A-relation}
        s_t=B_t-a_R^2A_t,
\end{equation}
and hence
\begin{equation}
\label{eq:B-s-A-relation}
        B_t=s_t+a_R^2A_t.
\end{equation}
Substituting this expression for \(B_t\) into \eqref{eq:robin-bridge-local} yields
\begin{equation}
\label{eq:robin-bridge-slope-form}
        x^R_{t,t+\tau}
        =
        X_t+
        \left[
        A_t+\left(s_t+a_R^2A_t\right)\tau
        \right]e^{-a_R^2\tau}.
\end{equation}
Equivalently,
\begin{equation}
\label{eq:robin-projection-form}
        x^R_{t,t+\tau}
        =
        \ell_R(\tau)^\top Z^R_t,
\end{equation}
where
\begin{equation}
\label{eq:ZR-robin-local}
        Z^R_t=
        \begin{pmatrix}
        A_t\\
        s_t\\
        X_t
        \end{pmatrix}
\end{equation}
and
\begin{equation}
\label{eq:ellR-robin-local}
        \ell_R(\tau)
        =
        \begin{pmatrix}
        (1+a_R^2\tau)e^{-a_R^2\tau}\\
        \tau e^{-a_R^2\tau}\\
        1
        \end{pmatrix}.
\end{equation}

This representation also clarifies the meaning of a matched-boundary state. In the Neumann/free bridge
\[
        x^N_{t,t+\tau}=X_t+(x^N_{t,t}-X_t)e^{-a^2\tau},
\]
the condition
\[
        x^N_{t,t}=X_t
\]
collapses the entire curve to the flat configuration \(x^N_{t,t+\tau}=X_t\). In the Robin bridge the analogous endpoint-matching condition is instead
\begin{equation}
\label{eq:robin-matched-boundary-A-zero}
        x^R_{t,t}=X_t
        \qquad\Longleftrightarrow\qquad
        A_t=0.
\end{equation}
This does not, by itself, collapse the Robin curve to a flat curve. Setting \(A_t=0\) in \eqref{eq:robin-bridge-slope-form} gives
\begin{equation}
\label{eq:robin-A-zero-interior-deformation}
        x^R_{t,t+\tau}
        =
        X_t+s_t\tau e^{-a_R^2\tau}.
\end{equation}
Thus the Robin endpoint can be matched to the long boundary while the interior still carries a repeated-root slope deformation. The fully flat matched state is the stronger condition \(A_t=0\) and \(s_t=0\). This is why the Robin boundary layer requires both the displacement coordinate \(A_t\) and the slope coordinate \(s_t\): the former measures endpoint mismatch, while the latter carries residual interior deformation after endpoint matching.

The vector \(Z^R_t=(A_t,s_t,X_t)^\top\) is a local coordinate representation of the Robin boundary layer. By a local coordinate representation we mean a finite-dimensional set of variables used to describe the boundary field over a fixed local regime. The Neumann representation uses \((x^N_{t,t},X_t)\). The Robin representation uses \((A_t,s_t,X_t)\). These coordinates describe different boundary geometries, so correlations measured in one representation should not be identified automatically with correlations measured in the other.

Suppose locally that the Robin boundary-layer coordinates satisfy
\begin{equation}
\label{eq:ZR-local-SDE}
        dZ^R_t=\mu_R(t)\,dt+\Gamma_R(t)\,dW_t,
\end{equation}
with instantaneous covariance matrix
\begin{equation}
\label{eq:SigmaR-local}
        \Sigma_R(t)=\Gamma_R(t)\Gamma_R(t)^\top.
\end{equation}
If \(a_R\) is fixed within the local regime, maturity dependence enters only through the deterministic loading vector \(\ell_R(\tau)\). Hence
\begin{equation}
\label{eq:dxR-projection-local}
        dx^R_{t,t+\tau}
        =
        \ell_R(\tau)^\top dZ^R_t,
\end{equation}
and the instantaneous variance at maturity \(\tau\) is
\begin{equation}
\label{eq:vR-projected-local}
        v_R(\tau,t)
        =
        \ell_R(\tau)^\top\Sigma_R(t)\ell_R(\tau).
\end{equation}
In applications one may approximate \(\Sigma_R(t)\) by a constant matrix within a fitted local sub-period. This is a local reduced-form approximation. It is not a direct consequence of the original two-state Neumann KFE.

This projected covariance is the natural Robin analogue of the two-state Neumann covariance. The structural Neumann correlation
\begin{equation}
\label{eq:rhoN-free-correlation}
        \rho_N=\operatorname{Corr}(dx^N_{t,t},dX_t)
\end{equation}
belongs to the free representation. In the Robin representation, the observable short endpoint is
\begin{equation}
        x^R_{t,t}=X_t+A_t,
\end{equation}
so correlations such as
\begin{equation}
        \operatorname{Corr}(dx^R_{t,t},dX_t)
\end{equation}
or
\begin{equation}
        \operatorname{Corr}(dA_t,dX_t)
\end{equation}
are transformed-coordinate correlations. They are informative about the Robin boundary layer, but they are not direct estimates of the primitive Neumann KFE correlation.

The scalar two-channel Robin MSNR formula is recovered only as a reduced case of this projected covariance framework. If the Robin representation is collapsed to two symmetric affine channels with common signal coefficient \(a_R\) and effective correlation \(\rho_R\), then
\begin{equation}
\label{eq:robin-msnr-reduced-recovered}
        \mathrm{MSNR}_{R}^{2}
        =
        \frac{2a_R^2}{1+\rho_R},
        \qquad
        \mathrm{MSNR}_{R}
        =
        |a_R|\sqrt{\frac{2}{1+\rho_R}}.
\end{equation}
The coefficient \(\rho_R\) in the scalar formula should therefore be interpreted
carefully.  It is not the primitive Neumann correlation
\[
        \rho_N=\operatorname{Corr}(dx^N_{t,t},dX_t),
\]
nor is it automatically equal to either of the directly observed Robin
diagnostic correlations
\[
        \operatorname{Corr}(dA_t,dX_t)
        \quad\text{or}\quad
        \operatorname{Corr}(dx^R_{t,t},dX_t).
\]
Rather, \(\rho_R\) denotes the effective correlation of the reduced
two-channel Robin projection obtained when the full three-coordinate
covariance
\[
        v_R(\tau,t)=\ell_R(\tau)^\top\Sigma_R(t)\ell_R(\tau)
\]
is collapsed to the symmetric scalar MSNR form.  In this sense \(\rho_R\)
is a maturity-invariant reduced coefficient only within the scalar Robin
special case; the general Robin boundary-layer representation is governed
by the full projected covariance matrix \(\Sigma_R(t)\).

There is therefore no canonical closed-form map
\[
        \Sigma_R(t) \longmapsto \rho_R(t)
\]
for a general three-coordinate Robin covariance matrix.  To define a scalar \(\rho_R\), one must also specify the two-channel reduction being used.  For example, on a chosen maturity grid \(\mathcal T\), one may define \(\rho_R\) operationally as the coefficient in the symmetric two-channel template that best approximates the projected variance curve \(\ell_R(\tau)^\top\Sigma_R(t)\ell_R(\tau)\).  Equivalently, if a scalar local MSNR level \(m_R(t)\) and a symmetric clock \(a_R(t)\) have already been fitted, the reduced coefficient is the value satisfying
\begin{equation}
\label{eq:rhoR-operational-definition}
        \rho_R(t)
        =
        \frac{2a_R(t)^2}{m_R(t)^2}-1.
\end{equation}
This definition is deliberately reduction-dependent.  The rank-two closure below constrains the three pairwise Robin correlations, but it does not by itself select a unique scalar \(\rho_R\) unless the reduction basis, weighting, and scalar clock are also specified.

\subsection{What \texorpdfstring{$\rho_R$}{rhoR} is, and is not}

The scalar coefficient \(\rho_R\) should be read operationally. It is the coefficient obtained when the full Robin projected covariance is fitted or projected into a symmetric two-channel MSNR form. It is therefore a reduced-form shadow of the three-coordinate covariance geometry, not a primitive boundary correlation.

In general,
\begin{equation}
\label{eq:rhoR-not-identities}
\rho_R\neq \rho_N,
\qquad
\rho_R\neq \operatorname{Corr}(dA_t,dX_t),
\qquad
\rho_R\neq \operatorname{Corr}(dx^R_{t,t},dX_t).
\end{equation}
These inequalities are not additional assumptions. They are reminders that the objects live in different coordinate systems. The free Neumann correlation \(\rho_N\) belongs to the structural two-state KFE. The Robin diagnostic correlations belong to the transformed boundary layer. The scalar \(\rho_R\) belongs only to the symmetric reduced MSNR approximation.

\begin{center}
\begin{tabular}{p{0.28\textwidth}p{0.47\textwidth}p{0.16\textwidth}}
\hline
Object & Interpretation & Status \\
\hline
\(\rho_N\) & Brownian correlation between the free Neumann short boundary and the stochastic long boundary & Structural \\
\(\operatorname{Corr}(dA_t,dX_t)\) & Robin amplitude--long-boundary correlation in the local boundary layer & Diagnostic \\
\(\operatorname{Corr}(dx^R_{t,t},dX_t)\) & Observable Robin endpoint--long-boundary correlation & Diagnostic \\
\(\rho_R\) & Effective scalar correlation after reducing the Robin projection to a symmetric two-channel MSNR form & Reduced-form \\
\hline
\end{tabular}
\end{center}

The coordinate hierarchy can be summarised as
\[
\begin{gathered}
\text{free Neumann coordinates}
\longrightarrow
\text{Robin boundary-layer coordinates} \\
\longrightarrow
\text{projected maturity covariance}
\longrightarrow
\text{scalar MSNR reduction}.
\end{gathered}
\]
At each arrow information may be transformed or compressed. The last object is useful when a compact scalar MSNR is desired, but the general Robin theory is the projected covariance theory \(\ell_R(\tau)^\top\Sigma_R(t)\ell_R(\tau)\).

The boundary-condition interpretation is as follows. In the continuous field formulation of Paper 2, a Robin boundary condition relates the boundary value to its normal derivative. In the finite-dimensional representation above, this motivates using both the short-end value \(x^R_{t,t}\) and the short-end slope \(s_t\) as boundary-layer coordinates. Empirically, this leads naturally to estimating local slope--level anchoring relations within sub-periods, but those estimation details belong to the empirical companion work rather than to the theoretical construction here.

\section{Rank-two Robin covariance closure}
\label{sec:rank-two-robin-covariance}

The projected Robin representation in Section 7 does not, by itself, impose a stochastic rank.  In the most general local representation the covariance matrix
\[
        \Sigma_R(t)=\operatorname{Cov}_t(dA_t,ds_t,dX_t)
\]
may have rank one, rank two, or rank three.  A rank-three specification treats the amplitude, slope, and long boundary as three independent local covariance directions.  A rank-one specification is the exact scalar MSNR limit in which all maturity-projected shocks are multiples of a single risk direction.  Between these extremes lies a particularly important local closure: the three Robin boundary-layer coordinates are allowed to span two stochastic directions.

This section records that rank-two closure.  It should be read as an additional local stochastic restriction, motivated by the affine return-to-risk discipline, rather than as a consequence of the Robin loading vector alone.  The purpose is to identify the covariance geometry that is preserved when the Robin boundary layer is organised by two independent boundary shocks.

Write
\begin{equation}
\label{eq:SigmaR-rank-two-general}
        \Sigma_R=
        \begin{pmatrix}
        \sigma_A^2 & \sigma_{sA} & \sigma_{AX}\\
        \sigma_{sA} & \sigma_s^2 & \sigma_{sX}\\
        \sigma_{AX} & \sigma_{sX} & \sigma_X^2
        \end{pmatrix},
\end{equation}
where the ordering of the coordinates is \((A_t,s_t,X_t)\).  The corresponding correlation matrix is
\begin{equation}
\label{eq:robin-correlation-matrix}
        R_R=
        \begin{pmatrix}
        1 & \rho_{sA} & \rho_{AX}\\
        \rho_{sA} & 1 & \rho_{sX}\\
        \rho_{AX} & \rho_{sX} & 1
        \end{pmatrix}.
\end{equation}
The three correlations in \eqref{eq:robin-correlation-matrix} are all Robin boundary-layer correlations.  They should not be identified with the structural Neumann correlation \(\rho_N\).

\subsection{Rank-two determinant condition}

If the local Robin covariance matrix has rank two and each marginal variance is non-zero, then the correlation matrix has determinant zero:
\begin{equation}
\label{eq:rank-two-corr-det-zero}
        \det R_R=0.
\end{equation}
Expanding gives
\begin{equation}
\label{eq:rank-two-corr-det-expanded}
        1+2\rho_{sA}\rho_{AX}\rho_{sX}
        -\rho_{sA}^2-\rho_{AX}^2-\rho_{sX}^2=0.
\end{equation}
Solving this quadratic for \(\rho_{sA}\) yields the boundary-correlation triangle
\begin{equation}
\label{eq:rank-two-corr-triangle}
\boxed{
        \rho_{sA}
        =
        \rho_{AX}\rho_{sX}
        \pm
        \sqrt{(1-\rho_{AX}^2)(1-\rho_{sX}^2)}.
}
\end{equation}
Thus, under rank-two closure, knowledge of the two correlations \(\rho_{AX}\) and \(\rho_{sX}\) determines the third correlation \(\rho_{sA}\), up to an orientation branch.

The same statement can be written directly in covariance units.  Assume \(\sigma_X^2>0\), and treat \(\sigma_{sA}\) as the unknown entry after observing
\[
        \sigma_A^2,
        \quad \sigma_s^2,
        \quad \sigma_X^2,
        \quad \sigma_{AX},
        \quad \sigma_{sX}.
\]
The rank-two condition \(\det\Sigma_R=0\) gives
\begin{equation}
\label{eq:rank-two-sigma-sA-branches}
\boxed{
        \sigma_{sA}
        =
        \frac{
        \sigma_{AX}\sigma_{sX}
        \pm
        \sqrt{
        (\sigma_A^2\sigma_X^2-\sigma_{AX}^2)
        (\sigma_s^2\sigma_X^2-\sigma_{sX}^2)
        }}{
        \sigma_X^2
        }.
}
\end{equation}
Equivalently, \(\sigma_{sA}=\rho_{sA}\sigma_s\sigma_A\) after applying \eqref{eq:rank-two-corr-triangle}.

Equation \eqref{eq:rank-two-corr-triangle} has a simple geometric interpretation.  In a rank-two covariance system, the standardised innovations of \(A_t\), \(s_t\), and \(X_t\) can be represented as three unit vectors in a plane.  Knowing \(\rho_{AX}\) fixes the angle between the amplitude and long-boundary vectors.  Knowing \(\rho_{sX}\) fixes the angle between the slope and long-boundary vectors.  The slope vector may lie on either side of the long-boundary vector, giving the two orientation branches in \eqref{eq:rank-two-corr-triangle}.  Once the branch is fixed locally, \(\rho_{sA}\) is no longer free.

\subsection{Relation to local MSNR discipline}

The exact maturity-invariant projected MSNR condition in a local Robin regime can be written as
\begin{equation}
\label{eq:local-maturity-invariant-msnr-rank-two-section}
        \frac{
        \ell_R(\tau)^\top\mu_R(t)
        }{
        \left(\ell_R(\tau)^\top\Sigma_R(t)\ell_R(\tau)\right)^{1/2}
        }
        =m_t,
        \qquad \tau\geq0.
\end{equation}
Here the common value \(m_t\) may vary over calendar time or across local regimes.  The invariance is across maturity, not across all dates.  In this sense MSNR is a local return-to-risk discipline on the maturity projection.

Condition \eqref{eq:local-maturity-invariant-msnr-rank-two-section} is stronger than rank-two covariance closure.  In its exact form it also restricts the local drift vector \(\mu_R(t)\).  The rank-two condition instead isolates the stochastic part of the same organising idea: the three observed Robin boundary-layer coordinates need not carry three independent covariance directions.  A local affine-preserving boundary field may therefore be studied through the pair
\[
        \text{MSNR discipline on projected return-to-risk,}
        \qquad
        \text{rank-two closure of projected covariance.}
\]
The former controls the drift--covariance alignment; the latter gives the directly testable covariance manifold.

A useful intermediate formulation is
\begin{equation}
\label{eq:mu-sigma-lambda-rank-two-section}
        \mu_R(t)=\Sigma_R(t)\lambda_R(t),
\end{equation}
where \(\lambda_R(t)\) is a local boundary price-of-risk, or signal, vector.  This equation allows the common signal structure to pass through the covariance matrix without requiring every realised maturity return-to-volatility ratio to be estimated accurately from noisy drift data.  Empirically, covariance restrictions such as \eqref{eq:rank-two-corr-triangle} are therefore cleaner diagnostics of the field geometry than direct estimates of \(\mu_R(t)\).

\subsection{Boundary-layer interpretation}

The rank-two closure says that there exists a non-zero vector \(q_t=(q_A,q_s,q_X)^\top\) such that, locally,
\begin{equation}
\label{eq:rank-two-linear-relation}
        q_A\,dA_t+q_s\,ds_t+q_X\,dX_t=0
\end{equation}
in covariance space.  Equivalently, one coordinate is locally covariance-spanned by the other two.  If \(q_s\neq0\), then
\begin{equation}
\label{eq:ds-spanned-by-A-X}
        ds_t=b_A(t)dA_t+b_X(t)dX_t
\end{equation}
for local coefficients \(b_A(t)\) and \(b_X(t)\).  This should not be confused with a deterministic pathwise identity.  It is a statement about the rank of the local innovation covariance.

A Robin boundary condition in the continuous field relates a boundary value to its normal derivative.  In the finite-dimensional boundary-layer representation, this motivates looking for local covariance-spanning relations involving the displacement \(A_t\), the slope \(s_t\), and the long boundary \(X_t\).  The most restrictive illustrative case is
\begin{equation}
\label{eq:robin-slope-increment-restrictive}
        ds_t\approx -\chi_t(dA_t+dX_t),
\end{equation}
which corresponds to a local slope--level anchoring relation with fixed \(\chi_t\).  This special case implies particular signs and magnitudes for \(\sigma_{sA}\) and \(\sigma_{sX}\).  The more general rank-two closure \eqref{eq:rank-two-corr-triangle} is weaker: it does not impose a specific anchoring coefficient, but it still requires the three Robin correlations to lie on the two-dimensional covariance manifold.

The theoretical role of the rank-two closure is therefore not to replace the full projected covariance matrix.  Rather, it identifies a parsimonious stochastic fabric within the Robin boundary layer.  When the closure holds locally, observing the amplitude--long and slope--long correlations determines the slope--amplitude correlation up to orientation.  This provides a sharp empirical implication of the boundary-coordinate construction while preserving the distinction between structural Neumann correlations, Robin diagnostic correlations, and reduced scalar MSNR coefficients.


\section{Symmetric two-boundary equilibrium density}
\label{sec:symmetric-equilibrium-density}

The preceding MSNR calculation is a local covariance statement. It is useful also to record the equilibrium density naturally associated with the symmetric affine-preserving closure. This connects the two-boundary construction back to the shifted inverse-gamma equilibrium law of the one-boundary IRFR model. The result in this section is not a classification of all two-boundary equilibria. It is the zero-flux density corresponding to the symmetric inverse-square-root covariance representation with common structural clock and constant correlation.

Thus the density below should be interpreted as the stationary law naturally paired with the compact symmetric covariance closure. It is not claimed to be the only possible two-boundary stationary law. Its role is to show that the same covariance geometry that produces the symmetric MSNR also generates a coherent equilibrium concentration kernel.

Let the positive residual boundary coordinates be denoted by
\[
x>0,
\qquad
X>0.
\]
In this section, \(x\) denotes the short residual coordinate of the symmetric two-boundary equilibrium law, not the Robin amplitude \(A_t\). The matched-boundary diagonal \(x=X\) therefore belongs to the two-boundary density state space. Its Robin analogue, in the boundary-layer representation of Section 7, is the endpoint-matching condition \(A_t=0\). Because Robin also carries the slope coordinate \(s_t\), this analogue need not be a flat-curve condition unless \(s_t=0\) as well.

In the shifted case these should be read as residuals from the corresponding equilibrium anchors. Thus, for raw states \(\hat x\) and \(\hat X\), one may write
\[
        x=\hat x-x_e,
        \qquad
        X=\hat X-X_e,
\]
with \(x,X>0\). The precise anchor values depend on the boundary normalisation and are not needed for the non-degenerate density below. Under the symmetric parameterisation the two boundary rate parameters are identified with the same clock, written here as \(a^2\), and the boundary correlation is a constant \(\rho\) with \(|\rho|<1\). Define
\[
p=\frac{\nu+3}{2},
\qquad
\lambda=\frac{a^2}{1-\rho^2}.
\]
The parameter \(\nu\) controls the power-law prefactor and boundary-tail behaviour of the equilibrium density. It does not enter the diagonal MSNR identity below, which is determined by the exponential concentration kernel. On the matched-boundary diagonal the prefactor contributes only the polynomial envelope \(x^{-2p}\); it does not change the exponential concentration rate, which is the object compared with the symmetric MSNR.

The symmetric two-boundary equilibrium density is
\begin{equation}
\label{eq:symmetric-two-boundary-density}
\boxed{
\phi(x,X)
=
C(xX)^{-p}
\exp\left[
-\lambda
\left(
\frac1x+
\frac1X-
\frac{2\rho}{\sqrt{xX}}
\right)
\right]
},
\qquad x>0,\ X>0 .
\end{equation}
Equivalently,
\begin{equation}
\label{eq:symmetric-two-boundary-density-expanded}
\phi(x,X)
=
C(xX)^{-(\nu+3)/2}
\exp\left[
-\frac{a^2}{1-\rho^2}
\left(
\frac1x+
\frac1X-
\frac{2\rho}{\sqrt{xX}}
\right)
\right].
\end{equation}
The case \(\rho=0\) should be interpreted carefully. Within the equilibrium family \eqref{eq:symmetric-two-boundary-density}, setting \(\rho=0\) removes the cross-boundary coupling term and gives
\begin{equation}
\label{eq:rho-zero-factorisation}
\phi(x,X)
=
C x^{-p}e^{-a^2/x}\,X^{-p}e^{-a^2/X}.
\end{equation}
After absorbing constants into the marginal normalisations, this is precisely the product of two one-boundary shifted inverse-gamma equilibrium densities with common rate \(a^2\):
\begin{equation}
\label{eq:rho-zero-independent-product}
\phi(x,X)=\phi_x(x)\phi_X(X).
\end{equation}
Thus, in this specified two-boundary equilibrium family, the uncorrelated closure \(\rho=0\) is the separable-density closure, which permits the stronger equilibrium interpretation that the two residual boundary channels are independent random variables. This is stronger than the generic probabilistic statement ``zero correlation implies independence,'' which is false for non-Gaussian variables in general. The implication here is model-specific: the coefficient \(\rho\) appears directly as the cross-coupling parameter in the equilibrium density. For non-zero \(\rho\), the cross term couples the two boundary residuals and the density no longer factorises.

The structure is clearest in inverse-square-root coordinates
\begin{equation}
\label{eq:inverse-root-coordinates}
u=x^{-1/2},
\qquad
v=X^{-1/2}.
\end{equation}
Since \(dx\,dX=4u^{-3}v^{-3}du\,dv\), the density transforms as
\begin{equation}
\label{eq:uv-density}
\phi(x,X)\,dx\,dX
=
4C u^{\nu}v^{\nu}
\exp\left[
-\frac{a^2}{1-\rho^2}
\left(u^2+v^2-2\rho uv\right)
\right]du\,dv .
\end{equation}
Thus the two-boundary density is a correlated inverse-square-root extension of the one-boundary shifted inverse-gamma law. It is normalisable for
\[
a^2>0,
\qquad
|\rho|<1,
\qquad
\nu>-1.
\]
The lower boundaries \(x=0\) and \(X=0\) are exponentially suppressed, while the upper tails are integrable precisely when \(\nu>-1\).

The restriction \(|\rho|<1\) gives the non-degenerate two-dimensional density. The singular endpoints are limiting cases rather than normalisable two-dimensional densities. At \(\rho=1\), the two inverse-square-root coordinates are perfectly comoving. In centred or standardised coordinates one may write
\begin{equation}
\label{eq:rho-one-collapse}
        \widetilde v=\widetilde u,
        \qquad
        \zeta=\widetilde u+\widetilde v=2\widetilde u .
\end{equation}
More generally, before standardisation, perfect positive correlation means that the two stochastic channels lie on a single affine line, so one may write \(V=\alpha U+\beta\) almost surely when the relevant second moments exist. In a particular matched-boundary normalisation one may choose anchors such as \(X_e=x_\infty\) and \(x_e=x_\infty/2\), but this belongs to the singular matched-boundary limit rather than to the non-degenerate \(|\rho|<1\) density. Thus \(\rho=1\) is the single-random-variable limit of the two-boundary representation, not a genuinely two-dimensional equilibrium density. Similarly, \(\rho=-1\) is the perfectly offsetting singular limit.

The zero-flux structure can be verified directly in the \((u,v)\) representation. Let
\[
\Sigma_{\rho}=
\begin{pmatrix}
1 & \rho\\
\rho & 1
\end{pmatrix},
\]
and let \(q(u,v)\) denote the density in \((u,v)\)-coordinates, omitting the normalising constant:
\[
q(u,v)
\propto
u^{\nu}v^{\nu}
\exp\left[
-\frac{a^2}{1-\rho^2}
\left(u^2+v^2-2\rho uv\right)
\right].
\]
For the reversible diffusion with constant covariance \(\Sigma_\rho\), zero flux requires
\begin{equation}
\label{eq:zero-flux-drift}
b=\frac12\Sigma_\rho\nabla\log q .
\end{equation}
The score derivatives are
\[
\partial_u\log q
=
\frac{\nu}{u}
-
\frac{2a^2}{1-\rho^2}(u-\rho v),
\qquad
\partial_v\log q
=
\frac{\nu}{v}
-
\frac{2a^2}{1-\rho^2}(v-\rho u).
\]
Substitution into \eqref{eq:zero-flux-drift} gives
\begin{align}
\label{eq:zero-flux-bu}
b_u
&=
-a^2u
+
\frac{\nu}{2}\left(\frac1u+\frac{\rho}{v}\right),\\
\label{eq:zero-flux-bv}
b_v
&=
-a^2v
+
\frac{\nu}{2}\left(\frac{\rho}{u}+\frac1v\right).
\end{align}
Hence \eqref{eq:symmetric-two-boundary-density} is the zero-flux equilibrium density for the symmetric inverse-square-root covariance closure. In the special case \(\nu=0\), the drift in \((u,v)\) reduces to the linear restoring form
\[
b_u=-a^2u,
\qquad
b_v=-a^2v,
\]
with all boundary repulsion encoded in the coordinate transformation back to \((x,X)\).

The link with the symmetric MSNR is obtained by restricting the exponential kernel of the density to the matched-boundary diagonal \(x=X\). This is a state-space restriction: it selects the boundary configurations in which the short and long residual coordinates coincide. It is not the same operation as taking a long-maturity limit. On this diagonal,
\[
\frac1x+
\frac1X-
\frac{2\rho}{\sqrt{xX}}
=
\frac{2(1-\rho)}{x}.
\]
Therefore the exponential rate becomes
\begin{equation}
\label{eq:diagonal-density-msnr-rate}
\frac{a^2}{1-\rho^2}\frac{2(1-\rho)}{x}
=
\frac1x\frac{2a^2}{1+\rho}.
\end{equation}
The diagonal exponential concentration coefficient is exactly the symmetric two-channel MSNR from \eqref{eq:symmetric-MSNR}:
\begin{equation}
\label{eq:density-msnr-identity}
\boxed{
\mathrm{MSNR}^2=\frac{2a^2}{1+\rho}.
}
\end{equation}
Thus, in the symmetric affine-preserving closure, the same structural clock that governs boundary transport and the two-channel covariance metric also governs the exponential concentration of the matched-boundary state. This statement concerns the exponential kernel of the density, not the full density itself; the algebraic prefactor \((xX)^{-p}\) remains part of the equilibrium law.

It is also useful to separate the common-mode concentration from the boundary-separation penalty. Since
\[
u^2+v^2-2\rho uv=(u-v)^2+2(1-\rho)uv,
\]
the exponent in \eqref{eq:symmetric-two-boundary-density} can be written as
\begin{equation}
\label{eq:separation-penalty-density}
-\frac{a^2}{1-\rho^2}
\left(
\frac{\sqrt{X}-\sqrt{x}}{\sqrt{xX}}
\right)^2
-
\frac{2a^2}{(1+\rho)\sqrt{xX}}.
\end{equation}
The first term penalises separation between the two boundary residuals in inverse-square-root geometry. The second term is the common-mode equilibrium concentration term. On the diagonal the separation penalty vanishes, leaving precisely the MSNR rate in the exponential kernel.

\section{Asymptotic boundary dynamics and long-end variance}

Taking the long-maturity limit of the minimal one-risk-direction boundary law gives
\begin{equation}
\label{eq:long-boundary-sde-X}
dX_t
=
-a^2(X_t-x_\infty)\,dt
+
\sqrt{2}a(X_t-x_\infty)\,dW_t^X.
\end{equation}
Equivalently, for \(Y_t=X_t-x_\infty\),
\begin{equation}
\label{eq:multiplicative-displacement-law}
dY_t
=
-a^2Y_t\,dt
+
\sqrt{2}aY_t\,dW_t^X.
\end{equation}
For \(Y_{t_0}\neq0\),
\begin{equation}
\label{eq:Y-explicit-solution}
Y_t
=
Y_{t_0}
\exp\left(-2a^2(t-t_0)+\sqrt{2}a(W_t^X-W_{t_0}^X)\right).
\end{equation}
The sign of the displacement is preserved. If \(Y_{t_0}=0\), then \(Y_t=0\) thereafter in this minimal multiplicative subclass.

The generator for the minimal long-boundary process is
\begin{equation}
\label{eq:long-boundary-generator}
(\mathcal A\phi)(X)
=
-a^2(X-x_\infty)\phi_X(X)
+
a^2(X-x_\infty)^2\phi_{XX}(X).
\end{equation}
Applying the generator to \(\phi(X)=X\) gives the drift \(-a^2(X-x_\infty)\). Applying it to \((X-X_0)^2\) at \(X=X_0\) identifies the variance coefficient
\[
\sigma_X^2(X)=2a^2(X-x_\infty)^2.
\]
Thus the generator corroborates the asymptotic SDE.

For the Robin implementation, the corresponding long-boundary variance limit is obtained by replacing the generic long-boundary clock \(a\) with the Robin clock \(a_R\). As \(\tau\to\infty\),
\[
A_R(t,\tau)\to0,
\qquad
C_R(t,\tau)\to Y_t.
\]
Therefore the Robin long-end variance limit is
\begin{equation}
\label{eq:long-end-variance-floor}
\lim_{\tau\to\infty}v_{\rho_R,R}(t,\tau)
=
2a_R^2Y_t^2
=
2a_R^2(X_t-x_\infty)^2.
\end{equation}
The cross term does not contribute to the asymptotic limit because the short/interior Robin channel dies out at the long boundary. Correlation matters at finite maturities, where both channels are active. The term ``limit is used deliberately: this quantity may vanish when the realised long-boundary state coincides with the equilibrium anchor, \(X_t=x_\infty\).


\section{Maturity-space Jacobian and stochastic rank}

The completed boundary-field representation separates maturity geometry from state dynamics.  It is useful to make explicit that there are two equivalent coordinate descriptions of the same maturity bridge.  The first is the canonical bridge representation, written in endpoint/slope coordinates.  The second is the Robin boundary-layer representation introduced in Section 7.

In canonical bridge coordinates, write
\begin{equation}
\label{eq:canonical-bridge-section10}
x_t(\tau)
=
x_{t,t}e^{-a^2\tau}
+
B_t\tau e^{-a^2\tau}
+
X_t\left(1-e^{-a^2\tau}\right).
\end{equation}
Equivalently,
\begin{equation}
\label{eq:canonical-jacobian-state-representation}
x_t(\tau)=J_c(\tau)S^c_t,
\end{equation}
where
\begin{equation}
\label{eq:boundary-jacobian}
J_c(\tau)
=
\left(
 e^{-a^2\tau},
 \tau e^{-a^2\tau},
 1-e^{-a^2\tau}
\right),
\end{equation}
and
\begin{equation}
\label{eq:canonical-state-vector}
S^c_t=
\begin{pmatrix}
x_{t,t}\\
B_t\\
X_t
\end{pmatrix}.
\end{equation}
The three entries of \(J_c(\tau)\) have distinct boundary meanings: short-boundary transport, interior slope deformation, and long-boundary transport.  This is the coordinate system in which the original bridge is most transparent.

The Robin boundary-layer coordinates are obtained by replacing the short endpoint and bridge slope with the transformed boundary variables
\begin{equation}
\label{eq:robin-coordinate-map-section10}
A_t=x^R_{t,t}-X_t,
\qquad
s_t=B_t-a_R^2A_t,
\qquad
Z^R_t=(A_t,s_t,X_t)^\top.
\end{equation}
In these coordinates, the same Robin maturity strip is written as
\begin{equation}
\label{eq:robin-loading-section10}
x^R_{t,t+\tau}=\ell_R(\tau)^\top Z^R_t,
\end{equation}
with
\begin{equation}
\label{eq:robin-loading-vector-section10}
\ell_R(\tau)
=
\begin{pmatrix}
(1+a_R^2\tau)e^{-a_R^2\tau}\\
\tau e^{-a_R^2\tau}\\
1
\end{pmatrix}.
\end{equation}
Thus \(J_c(\tau)\) and \(\ell_R(\tau)\) are not competing maturity geometries.  They are alternative coordinate descriptions of the same boundary-field bridge: \(J_c(\tau)\) uses the canonical variables \((x_{t,t},B_t,X_t)\), while \(\ell_R(\tau)\) uses the Robin boundary-layer variables \((A_t,s_t,X_t)\).

The canonical loadings are linearly independent.  Indeed, if
\[
c_1e^{-a^2\tau}+c_2\tau e^{-a^2\tau}+c_3(1-e^{-a^2\tau})=0
\]
for all \(\tau\ge0\), then letting \(\tau\to\infty\) gives \(c_3=0\).  The remaining identity
\[
(c_1+c_2\tau)e^{-a^2\tau}=0
\]
forces \(c_1=c_2=0\).  The same reasoning applies to the Robin loading family: the constant long-boundary loading, the exponentially decaying amplitude loading, and the exponentially decaying slope loading define distinct maturity directions.

Neither coordinate representation imposes a particular stochastic rank.  In canonical coordinates, one may write
\[
dS^c_t=\mu^c_t\,dt+Q^c_t\,dW_t,
\]
with state covariance \(\Sigma^c_t=Q^c_t(Q^c_t)^\top\).  The corresponding maturity covariance kernel is
\begin{equation}
\label{eq:jacobian-covariance-kernel}
K^c_t(\tau,\sigma)=J_c(\tau)\Sigma^c_tJ_c(\sigma)^\top.
\end{equation}
In Robin coordinates, the analogous representation is
\[
dZ^R_t=\mu^R_t\,dt+Q^R_t\,dW_t,
\qquad
\Sigma^R_t=Q^R_t(Q^R_t)^\top,
\]
and
\begin{equation}
\label{eq:robin-covariance-kernel-section10}
K^R_t(\tau,\sigma)=\ell_R(\tau)^\top\Sigma^R_t\ell_R(\sigma).
\end{equation}
A one-direction, two-direction, or higher-rank volatility model corresponds to different restrictions on the relevant state covariance matrix.  The invariant contribution of the boundary theory is not a pre-imposed stochastic rank, but the maturity geometry induced by the boundary coordinates.  The rank-two Robin closure of Section~\ref{sec:rank-two-robin-covariance} is therefore best read as a local stochastic specialisation of this geometry: it is a testable covariance restriction, not an assumption required by the loading vector itself.

\section{Relation to risk-neutral pricing}

The stochastic long-boundary extension is a physical-field construction, written under the physical measure \(\mathbb P\). It identifies how a realised long-maturity boundary state can move within the same boundary-field architecture and how drift must be paired with variance when short/interior and long-boundary shocks are not perfectly correlated.

The money-market account is generated by the realised short-end roll \(x_{t,t}\). The stochastic long-boundary \(X_t\) does not directly fund the bank account. Instead, it enters bond prices through the full IRFR strip
\[
\{x_t(u-t):u\in[t,T]\}.
\]
A risk-neutral version, written under \(\mathbb Q\), would require specifying the market price of boundary risk and imposing the drift restrictions needed for arbitrage-free bond pricing. A full no-arbitrage treatment must trace \(X_t\), the Robin interior state, and the boundary-risk covariance structure through the risk-neutral drift restriction and the rolling-maturity transport equation. A companion pricing treatment can impose risk-neutral dynamics on the same volatility-aligned geometry and derive the associated no-arbitrage bond-pricing equations, in the tradition of term-structure volatility and option-pricing models with maturity-dependent volatility structures \cite{RitchkenChuang}. The present paper does not impose that pricing restriction. It completes the physical boundary-field architecture and identifies the admissible long-boundary geometry. The compatibility of the current IRFR framework with the more conventional Heath-Jarrow-Morton (HJM) framework can be found in Appendix \ref{app:hjm-pricing}.


\section{Empirical Implications}

The coefficient-matching theorem identifies constant correlation as part of the clean affine admissible baseline in the free two-state representation. Thus empirical work should distinguish the structural free-coordinate correlation from correlations measured after imposing a Robin boundary transformation.

In the Robin branch the theory supplies the local projection geometry
\[
Z^R_t=(A_t,s_t,X_t)^\top,
\qquad
x^R_{t,t+\tau}=\ell_R(\tau)^\top Z^R_t,
\]
with
\[
\ell_R(\tau)=
\left(
(1+a_R^2\tau)e^{-a_R^2\tau},
\tau e^{-a_R^2\tau},
1
\right)^\top.
\]
The maturity variance is then estimated through
\[
        v_R(\tau)=\ell_R(\tau)^\top\Sigma_R\ell_R(\tau),
\]
where \(\Sigma_R\) is the local covariance matrix of \((A_t,s_t,X_t)\), approximated empirically within a chosen Robin sub-period.

This provides a more precise empirical target than a single realised return-to-volatility ratio. The scalar expression
\[
        \mathrm{MSNR}_{R}^{2}=\frac{2a_R^2}{1+\rho_R}
\]
remains useful as a symmetric two-channel reduction, but empirical correlations such as \(\operatorname{Corr}(dA_t,dX_t)\) or \(\operatorname{Corr}(dx^R_{t,t},dX_t)\) should be treated as Robin boundary-layer diagnostics rather than direct measurements of the primitive Neumann correlation.

The rank-two closure gives an additional empirical target.  In a local Robin regime, define
\[
        \rho_{AX}=\operatorname{Corr}(dA_t,dX_t),
        \qquad
        \rho_{sX}=\operatorname{Corr}(ds_t,dX_t),
        \qquad
        \rho_{sA}=\operatorname{Corr}(ds_t,dA_t).
\]
If the boundary-layer covariance is locally rank two, then
\[
        \rho_{sA}
        =
        \rho_{AX}\rho_{sX}
        \pm
        \sqrt{(1-\rho_{AX}^2)(1-\rho_{sX}^2)}.
\]
Thus the slope--amplitude correlation is not an independent object once the two long-boundary correlations and the orientation branch are fixed.  This is a covariance-level implication of the boundary geometry and is more stable than direct drift-based MSNR estimation.

The stronger empirical test is whether the completed Robin projection geometry improves the joint description of curve shape, maturity volatility, and boundary-layer covariance across local sub-periods, and whether the rank-two boundary-correlation triangle continues to hold in blind time windows that are not selected by the Robin segmentation.


\section{Summary}

The stochastic long-boundary extension separates the equilibrium anchor \(x_\infty\) from the realised long-boundary state \(X_t\). The displacement \(Y_t=X_t-x_\infty\) is transported through maturity by \(1-e^{-a^2\tau}\). Together with the short-boundary loading \(e^{-a^2\tau}\) and the Robin interior loading \(\tau e^{-a^2\tau}\), this gives the completed maturity-space Jacobian
\[
J(\tau)=
\left(
 e^{-a^2\tau},
 \tau e^{-a^2\tau},
 1-e^{-a^2\tau}
\right).
\]

Coefficient matching in the two-variable KFE shows that the long-boundary extension preserves an affine admissible subclass. The necessary geometric condition is volatility alignment,
\[
y_x=\frac{g_y}{g_x},
\qquad
Y_X=\frac{g_Y}{g_X}.
\]
On the orientation-preserving branch, the free two-boundary representation has a single structural Brownian correlation,
\[
\rho_{yY}=\rho_{xX}=\rho.
\]
Under affine \(g\), constant \(f/g\), and this constant-correlation branch, the transformed two-boundary system has a global correlated MSNR:
\[
\mathrm{MSNR}_{y,Y}^2
=
\frac{\alpha^2-2\rho\alpha\beta+\beta^2}{1-\rho^2}.
\]
In the symmetric case \(\alpha=\beta=a\), this reduces to
\[
\mathrm{MSNR}_{y,Y}^2=\frac{2a^2}{1+\rho}.
\]

The Robin branch is not obtained by simply replacing the free short coordinate with a fitted Robin endpoint. The free Neumann short endpoint \(x^N_{t,t}\), the pinned Dirichlet endpoint \(x_D\), and the partially anchored Robin endpoint \(x^R_{t,t}\) should be distinguished. The heuristic relation
\[
x^R_{t,t}\approx \omega_t x^N_{t,t}+(1-\omega_t)x_D
\]
serves as a roadmap for the Robin boundary-layer coordinates rather than as an imposed theorem. The local Robin representation is
\[
Z^R_t=(A_t,s_t,X_t)^\top,
\qquad
A_t=x^R_{t,t}-X_t,
\qquad
s_t=\left.\partial_\tau x^R_{t,t+\tau}\right|_{\tau=0}.
\]
The maturity strip is generated by
\[
x^R_{t,t+\tau}=\ell_R(\tau)^\top Z^R_t,
\qquad
\ell_R(\tau)=
\left(
(1+a_R^2\tau)e^{-a_R^2\tau},
\tau e^{-a_R^2\tau},
1
\right)^\top.
\]
Thus the general Robin maturity variance is the projected covariance
\[
v_R(\tau,t)=\ell_R(\tau)^\top\Sigma_R(t)\ell_R(\tau).
\]

The canonical maturity-space Jacobian and the Robin loading vector are therefore alternative coordinate descriptions of the same boundary-field bridge: the former uses \((x_{t,t},B_t,X_t)\), while the latter uses the boundary-layer variables \((A_t,s_t,X_t)\).

The local rank-two Robin covariance closure isolates the stochastic fabric associated with this boundary layer.  If \(\Sigma_R=\operatorname{Cov}(dA_t,ds_t,dX_t)\) has rank two and the marginal variances are non-zero, then the Robin correlation matrix satisfies
\[
        1+2\rho_{sA}\rho_{AX}\rho_{sX}
        -\rho_{sA}^2-\rho_{AX}^2-\rho_{sX}^2=0.
\]
Equivalently,
\[
        \rho_{sA}
        =
        \rho_{AX}\rho_{sX}
        \pm
        \sqrt{(1-\rho_{AX}^2)(1-\rho_{sX}^2)}.
\]
Thus, in the rank-two closure, the slope--amplitude boundary correlation is determined by the two long-boundary correlations up to orientation.  This covariance relation is the directly testable stochastic counterpart of the local MSNR return-to-risk discipline.

The scalar expression
\[
\mathrm{MSNR}_{R}^2=\frac{2a_R^2}{1+\rho_R}
\]
is recovered only as a symmetric two-channel reduction of this projected covariance framework. In the Neumann bridge, the matched condition \(x^N_{t,t}=X_t\) gives a flat curve. In the Robin bridge, the analogous endpoint condition is \(A_t=0\); the curve is flat only under the stronger condition \(A_t=0\) and \(s_t=0\).

The symmetric affine closure also admits a coherent two-boundary equilibrium density. The resulting density is a correlated inverse-square-root extension of the one-boundary shifted inverse-gamma law. Along the matched-boundary diagonal of the two-boundary density, the exponential concentration coefficient is exactly the symmetric MSNR. This connects the affine admissibility theorem, the two-boundary covariance geometry, and the equilibrium density hierarchy.

The empirical implication is that boundary-regime correlations are coordinate dependent. Correlations in the free two-boundary representation estimate the primitive KFE geometry. Correlations measured from Robin variables such as \(A_t\), \(s_t\), or \(x^R_{t,t}\) diagnose the transformed Robin boundary layer. The appropriate empirical target is therefore not only a scalar realised MSNR, but the joint performance of the completed boundary geometry, the local Robin projection covariance, and the associated maturity-volatility structure.



\clearpage
\begin{appendices}
\section{Compatibility with risk-neutral HJM pricing}
\label{app:hjm-pricing}

The main text develops the boundary-coordinate geometry of the instantaneous
forward-rate field under physical boundary dynamics.  This appendix records
how the same coordinate representation may be embedded in a standard
risk-neutral Heath--Jarrow--Morton representation.  The purpose is not to
replace the HJM no-arbitrage condition, but to clarify that the boundary
coordinates used in the paper are compatible with risk-neutral pricing. The physical-measure construction in the main text is not claimed to satisfy the HJM drift restriction under \(\mathbb P\); the restriction appears only after choosing an equivalent risk-neutral measure and imposing the corresponding no-arbitrage drift under \(\mathbb Q\).

Let \(f(t,T)\) denote the instantaneous forward rate for maturity \(T\), and
write \(\tau=T-t\).  Under a risk-neutral measure \(\mathbb Q\), a standard
HJM representation takes the form
\[
        df(t,T)
        =
        \alpha^{\mathbb Q}(t,T)\,dt
        +
        \sum_{k=1}^{m}\sigma_k(t,T)\,dW^{\mathbb Q}_{k,t}.
\]
In the absence of arbitrage, the drift is determined by the volatility
structure:
\[
        \alpha^{\mathbb Q}(t,T)
        =
        \sum_{k=1}^{m}
        \sigma_k(t,T)
        \int_t^T \sigma_k(t,u)\,du .
\]
Thus, once the risk-neutral volatility field is specified, the HJM drift
restriction supplies the corresponding no-arbitrage drift.

The boundary-coordinate representation used in the main text may be viewed
as a structured specification of the forward-rate field,
\[
        f(t,T)=\mathcal F(\tau;Z_t),
\]
where \(Z_t\) is a finite-dimensional vector of boundary coordinates.  In
the free Neumann representation one may take
\[
        Z^N_t=(x^N_{t,t},X_t)^\top,
\]
whereas in the Robin boundary-layer representation one may take
\[
        Z^R_t=(A_t,s_t,X_t)^\top .
\]
The choice of \(Z_t\) is a coordinate restriction on the curve; it is not a
substitute for the HJM drift condition.

Suppose that, under \(\mathbb Q\), the boundary coordinates satisfy
\[
        dZ_t
        =
        \mu_Z^{\mathbb Q}(t)\,dt
        +
        \Gamma_Z(t)\,dW^{\mathbb Q}_t,
\]
with instantaneous covariance matrix
\[
        \Sigma_Z(t)=\Gamma_Z(t)\Gamma_Z(t)^\top .
\]
Applying Itô's formula to \(f(t,T)=\mathcal F(\tau;Z_t)\) gives
\[
        df(t,T)
        =
        \left[
        -\partial_{\tau}\mathcal F(\tau;Z_t)
        +
        \nabla_Z\mathcal F(\tau;Z_t)^\top\mu_Z^{\mathbb Q}(t)
        +
        \frac12
        \operatorname{tr}
        \left(
        \Gamma_Z(t)\Gamma_Z(t)^\top
        \nabla_Z^2\mathcal F(\tau;Z_t)
        \right)
        \right]dt
        +
        \nabla_Z\mathcal F(\tau;Z_t)^\top
        \Gamma_Z(t)\,dW_t^{\mathbb Q}.
\]
Hence the induced HJM volatility is
\[
        \sigma_f(t,T)
        =
        \nabla_Z\mathcal F(\tau;Z_t)^\top\Gamma_Z(t).
\]
The HJM no-arbitrage drift is then obtained by imposing
\[
        \alpha^{\mathbb Q}(t,T)
        =
        \sigma_f(t,T)
        \int_t^T \sigma_f(t,u)^\top\,du,
\]
with the obvious vector generalisation when the Brownian dimension exceeds
one.

For the Robin bridge considered in the main text,
\[
        x^R_{t,t+\tau}
        =
        X_t+
        \left(A_t+B_t\tau\right)e^{-a_R^2\tau},
\]
where
\[
        s_t
        =
        \left.
        \partial_\tau x^R_{t,t+\tau}
        \right|_{\tau=0}
        =
        B_t-a_R^2A_t.
\]
Equivalently,
\[
        x^R_{t,t+\tau}
        =
        \ell_R(\tau)^\top Z^R_t,
        \qquad
        Z^R_t=
        \begin{pmatrix}
        A_t\\
        s_t\\
        X_t
        \end{pmatrix},
\]
with
\[
        \ell_R(\tau)
        =
        \begin{pmatrix}
        (1+a_R^2\tau)e^{-a_R^2\tau}\\
        \tau e^{-a_R^2\tau}\\
        1
        \end{pmatrix}.
\]
If
\[
        dZ^R_t
        =
        \mu_R^{\mathbb Q}(t)\,dt
        +
        \Gamma_R(t)\,dW_t^{\mathbb Q},
\]
then the induced risk-neutral forward-rate volatility is
\[
        \sigma_R(t,T)
        =
        \ell_R(\tau)^\top\Gamma_R(t),
\]
and the instantaneous variance is
\[
        \operatorname{Var}^{\mathbb Q}_t[df(t,T)]
        =
        \ell_R(\tau)^\top
        \Sigma_R(t)
        \ell_R(\tau)\,dt,
        \qquad
        \Sigma_R(t)=\Gamma_R(t)\Gamma_R(t)^\top .
\]
This is the projected covariance object used in the main text.

The HJM drift restriction then determines the risk-neutral drift associated
with this projected volatility:
\[
        \alpha_R^{\mathbb Q}(t,T)
        =
        \sigma_R(t,T)
        \int_t^T
        \sigma_R(t,u)^\top\,du .
\]
Thus the Robin boundary-layer geometry may be admitted into a standard
risk-neutral pricing framework by treating the boundary coordinates as the
state variables generating the forward-rate volatility field, while the HJM
condition supplies the corresponding no-arbitrage drift.

The distinction is important.  The boundary-coordinate analysis in the main
text is designed to expose the geometry of the curve, the Jacobian of
coordinate transformations, and the covariance quantities such as MSNR.  The
risk-neutral HJM representation supplies the no-arbitrage pricing
restriction.  These are complementary descriptions: the former structures
the volatility geometry, while the latter imposes the pricing drift under
\(\mathbb Q\).
\end{appendices}

\begin{thebibliography}{9}

\bibitem{OzaIRFR}
Oza, S. (2026).
\emph{An Axiomatically Consistent Single-Factor Term Structure Model with Admissible Humped Volatility}.
SSRN Working Paper, posted May 2026.
Available at SSRN: \url{https://papers.ssrn.com/sol3/papers.cfm?abstract_id=6454299}.

\bibitem{OzaShortBoundary}
Oza, S. (2026).
\emph{Unified Boundary Field Theory of Interest Rates II: Monetary Policy as a Short-End Boundary Condition}.
Zenodo, DOI {10.5281/zenodo.20943769},
\url{https://doi.org/10.5281/zenodo.20943769},
Preprint.

\bibitem{RitchkenChuang}
Ritchken, P., and Chuang, I. (2000).
Interest rate option pricing with volatility humps.
Review of Derivatives Research, 3, 237--262.

\end{thebibliography}
\end{document}
